\documentclass[12pt]{article}
\usepackage[utf8]{inputenc}
\usepackage[letterpaper,margin=1.0in]{geometry}
%% Some formatting stuff
\usepackage{authblk}
\usepackage{fancyhdr}
%\usepackage{lineno}
\usepackage{siunitx}
\usepackage{hyperref}
\usepackage{booktabs, array, longtable}
\usepackage{multirow}
\pagestyle{fancy}
\setlength{\headheight}{14.5pt} % Fix fancyhdr warning
% Custom citep command for biblatex
\newcommand{\citep}{\parencite}
\newcommand{\citet}{\textcite}
\fancyhead[R]{\textbf{pg\_gpu}}
% for figures
\usepackage{graphicx}
\usepackage{wrapfig}
\usepackage{float}
\usepackage{breakcites}
\usepackage{nameref}
\usepackage{amsmath}
\usepackage{xspace}

\graphicspath{ {./figures/} }

\newcommand{\comment}[1]{{\color{blue} #1}}
\newcommand{\pggpu}{\texttt{pg\_gpu}\xspace}

% for hyperlinks
\hypersetup{
    colorlinks=true,
    citecolor=black,
    urlcolor=cyan,
    linkcolor=blue
    }
\urlstyle{same}

% for highlighting text
\usepackage{xcolor}
\usepackage{soul}
\usepackage{tikz}

% code listings
\usepackage{listings}
\lstset{
    basicstyle=\ttfamily\small,
    breaklines=true,
    frame=single,
    language=Python
}

% bibliography
\usepackage[backend=biber,style=authoryear]{biblatex}
\addbibresource{refs.bib}

%\linenumbers
\renewcommand*{\bibfont}{\fontsize{10}{12}\selectfont}

\newcommand{\beginsupplement}{%
    \setcounter{table}{0}
    \renewcommand{\thetable}{S\arabic{table}}%
    \setcounter{figure}{0}
    \renewcommand{\thefigure}{S\arabic{figure}}%
    \renewcommand{\theHtable}{supp.\arabic{table}}%
    \renewcommand{\theHfigure}{supp.\arabic{figure}}%
}

\title{GPU accelerated population genetics statistics using \texttt{pg\_gpu}}
\author[1, 2]{TODO Authors}
\author[1, 2]{Nathaniel S. Pope}
\author[1, 2, *]{Andrew D. Kern}
\affil[1]{\small{University of Oregon, Institute of Ecology and Evolution}}
\affil[2]{\small{University of Oregon, Department of Biology}}
\affil[*]{\small{Corresponding author: \href{mailto:adkern@uoregon.edu}{adkern@uoregon.edu}}}

\date{\small{\today{}}}

\begin{document}

\maketitle

\begin{abstract}
TODO: Abstract goes here.
\end{abstract}

\section*{Introduction}

Population genetics relies on a rich toolkit of summary statistics to characterize
genetic variation within and between populations. Statistics such as nucleotide
diversity ($\pi$), Tajima's $D$, $F_{ST}$, and linkage disequilibrium (LD) are
foundational to studies of demographic history, natural selection, population
structure, and admixture. As whole-genome sequencing datasets have grown from
hundreds to hundreds of thousands of individuals---exemplified by projects such
as the Anopheles 1000 Genomes \citep{ag1000g2017} and UK Biobank
\citep{bycroft2018}---the computational cost of these analyses has become a
serious bottleneck.

A number of software tools have been developed to compute population genetics
statistics from genomic data. VCFtools \citep{danecek2011} was among the first
general-purpose tools, providing command-line access to diversity, $F_{ST}$,
and LD calculations from VCF files, but was designed for the sample sizes of
the early 1000 Genomes era. \texttt{scikit-allel} \citep{miles2024} brought
these analyses into the Python ecosystem with NumPy-based array operations,
offering a broader repertoire of statistics including haplotype-based selection
scans and $F$-statistics; however, it is no longer under active development.
\texttt{pixy} \citep{korunes2021} addressed the important problem of missing
data bias in estimates of $\pi$, $d_{XY}$, and $F_{ST}$, but is limited to
those three statistics. At the other end of the spectrum, PLINK
\citep{purcell2007, chang2015} and ANGSD \citep{korneliussen2014} are
high-performance C/C++ tools but are oriented toward GWAS and genotype
likelihood workflows, respectively, rather than the summary statistics used in
evolutionary population genetics. PopGenome \citep{pfeifer2014} offered a
comprehensive R-based toolkit but has been removed from CRAN and is no longer
maintained. For tree sequence data, tskit \citep{jeffery2026} provides
exceptionally efficient computation of a general class of statistics, but
requires data in tree sequence format rather than the VCF and Zarr formats
that remain the standard output of variant calling pipelines. Across all of
these tools, computation is performed on the CPU and statistics are typically
computed one at a time across genomic windows. For large-scale projects
involving windowed genome scans of multiple statistics across many populations,
wall-clock times of hours to days are common. This limitation is particularly
acute for two-locus statistics such as LD, where pairwise computations over
variant sites are inherently $O(n^2)$ and form a bottleneck in inference
pipelines such as \texttt{moments} \citep{jouganous2017}.

Modern graphics processing units (GPUs) offer massive parallelism that is
well-suited to the data-parallel structure of population genetics computation.
Many statistics reduce to operations over matrices of haplotypes or genotypes
that can be decomposed into independent per-site or per-window calculations,
making them natural candidates for GPU acceleration. While GPUs have been
applied to specific problems in population genetics---notably forward-in-time
simulation \citep{gofish2017} and demographic inference via diffusion
approximations \citep{gutenkunst2021}---no general-purpose GPU-accelerated
library for computing population genetics summary statistics from empirical
data has been available.

Here we present \pggpu, a Python library that leverages NVIDIA GPUs via CuPy
\citep{cupy2017} to compute a comprehensive catalog of population genetics
statistics. \pggpu implements over 70 functions spanning diversity, divergence,
the site frequency spectrum, linkage disequilibrium, haplotype-based selection
scans, admixture statistics, PCA, and relatedness---covering the large majority
of statistics in routine use. The library introduces fused CUDA kernels that
compute multiple statistics in a single pass over the data, and implements the
generalized theta framework of \citet{achaz2009} to efficiently derive all
standard diversity and neutrality test statistics from a single site frequency
spectrum computation. \pggpu provides a drop-in replacement for the LD parsing
module of \texttt{moments}, achieving up to 275-fold speedups while maintaining
numerical accuracy at machine precision. Across windowed genome scans, \pggpu
delivers up to 60-fold speedups relative to \texttt{scikit-allel}. \pggpu is
open source and freely available at
\url{https://github.com/kr-colab/pg_gpu}.

\section*{Design and Implementation}

TODO: Describe the design and implementation of \pggpu.

\section*{Features}

\pggpu provides a comprehensive catalog of population genetics statistics
organized into ten categories (Table~\ref{tab:features}). A complete listing of
all implemented statistics with citations is provided in the Supplementary
Material (Table~\ref{tab:full_catalog}).

\begin{table}[ht]
\centering
\small
\caption{Summary of statistics implemented in \pggpu. The full catalog with
citations is provided in Supplementary Table~\ref{tab:full_catalog}.}
\label{tab:features}
\begin{tabular}{llp{5.5cm}}
\toprule
\textbf{Category} & \textbf{Count} & \textbf{Representative statistics} \\
\midrule
Diversity \& neutrality tests & 20+ & $\pi$, $\theta_W$, $\theta_H$, Tajima's $D$, Fay \& Wu's $H$ \\
Divergence & 10+ & $F_{ST}$ (Hudson; Weir--Cockerham), $d_{XY}$, $D_a$, PBS \\
Linkage disequilibrium & 10+ & $r^2$, $D^2$, $D_z$, $\pi_2$, $Z_{nS}$, $\omega$ \\
Selection scans & 8 & iHS, nSL, XP-EHH, $H_{12}$, $H_{2}/H_{1}$ \\
Site frequency spectrum & 6+ & Unfolded, folded, scaled, joint (2-pop) SFS \\
Admixture / $F$-statistics & 6+ & Patterson's $F_2$, $F_3$, $D$ (ABBA-BABA) \\
Dimensionality reduction & 3 & PCA, randomized PCA, PCoA \\
Relatedness / kinship & 2 & GRM (GCTA-equivalent), IBS (PLINK-equivalent) \\
Distance distribution & 4 & Hamming distance, variance, skewness, kurtosis \\
Achaz framework & -- & Generalized $\theta$ estimators via SFS weight vectors \\
\bottomrule
\end{tabular}
\end{table}

TODO: Describe the features and statistics available in \pggpu.

\section*{Results}

\subsection*{Numerical accuracy}

TODO: Validation of \pggpu against established tools (scikit-allel, moments, PLINK)
using Ag1000G data. Report maximum relative error across statistic categories.

\subsection*{Performance benchmarks}

TODO: Wall-clock comparisons of \pggpu vs.\ scikit-allel for windowed genome scans
($\pi$, $\theta_W$, Tajima's $D$, $F_{ST}$, $d_{XY}$) across window sizes.
LD parsing benchmarks vs.\ moments. PCA/GRM benchmarks vs.\ PLINK.

\subsection*{Scaling}

TODO: Runtime scaling with number of samples and number of variants.
GPU memory usage at different dataset sizes.
Advantage of fused multi-statistic computation vs.\ serial single-statistic passes.

\subsection*{Application to Ag1000G}

TODO: End-to-end analysis on Ag1000G data demonstrating a complete workflow:
data loading, windowed statistics, selection scans, PCA, and total runtime.

\section*{Discussion}

TODO: Discussion goes here.

\section*{Availability}

\pggpu is freely available at \url{https://github.com/kr-colab/pg_gpu}.

\section*{Acknowledgments}

TODO: Acknowledgments go here.

\printbibliography

\clearpage
\beginsupplement

\section*{Supplementary Material}

\begin{longtable}{lll}
\caption{Complete catalog of statistics implemented in \pggpu.}
\label{tab:full_catalog} \\
\toprule
\textbf{Category} & \textbf{Statistic} & \textbf{Reference} \\
\midrule
\endfirsthead
\toprule
\textbf{Category} & \textbf{Statistic} & \textbf{Reference} \\
\midrule
\endhead
\midrule
\multicolumn{3}{r}{\textit{Continued on next page}} \\
\bottomrule
\endfoot
\bottomrule
\endlastfoot

% Diversity & neutrality tests
\multirow{12}{*}{Diversity} & Nucleotide diversity ($\pi$) & \citep{nei1979} \\
 & Watterson's $\theta_W$ & \citep{watterson1975} \\
 & Fay \& Wu's $\theta_H$ & \citep{fay2000} \\
 & Zeng's $\theta_L$ & \citep{zeng2006} \\
 & Tajima's $D$ & \citep{tajima1989} \\
 & Fay \& Wu's $H$, normalized $H^*$ & \citep{zeng2006} \\
 & Zeng's $E$ & \citep{zeng2006} \\
 & Segregating sites, singletons & -- \\
 & Haplotype diversity, count & -- \\
 & Expected / observed heterozygosity & -- \\
 & Wright's $F$ & \citep{wright1951} \\
 & Max DAF, DAF histogram & -- \\
\midrule

% Divergence
\multirow{10}{*}{Divergence} & $F_{ST}$ (Hudson) & \citep{hudson1992} \\
 & $F_{ST}$ (Weir--Cockerham) & \citep{weir1984} \\
 & $G_{ST}$ & \citep{nei1973} \\
 & $d_{XY}$ & \citep{nei1987} \\
 & $D_a$ (net divergence) & \citep{nei1979} \\
 & PBS & \citep{yi2010} \\
 & $S_{nn}$ & \citep{hudson2000} \\
 & $G_{\min}$ & \citep{geneva2015} \\
 & $d_d$, $d_d$ rank & \citep{schrider2018} \\
 & $Z_x$ & \citep{schrider2018} \\
\midrule

% Linkage disequilibrium
\multirow{8}{*}{LD} & $r$, $r^2$ & -- \\
 & $D^2$, $D_z$, $\pi_2$ & \citep{ragsdale2020} \\
 & $\sigma_d^2$ (unbiased) & \citep{ragsdale2020} \\
 & Kelly's $Z_{nS}$ & \citep{kelly1997} \\
 & Kim \& Nielsen's $\omega$ & \citep{kim2004} \\
 & $\mu_{\mathrm{LD}}$ (RAiSD) & \citep{alachiotis2018} \\
 & Fused windowed LD & -- \\
 & \texttt{moments} drop-in LD parsing & -- \\
\midrule

% Selection scans
\multirow{6}{*}{Selection} & iHS & \citep{voight2006} \\
 & nSL & \citep{ferrer-admetlla2014} \\
 & XP-EHH & \citep{sabeti2007} \\
 & XP-nSL & \citep{szpiech2021} \\
 & $H_1$, $H_{12}$, $H_{123}$, $H_2/H_1$ & \citep{garud2015} \\
 & EHH decay & \citep{sabeti2002} \\
\midrule

% SFS
\multirow{3}{*}{SFS} & Unfolded / folded SFS & -- \\
 & Scaled SFS & -- \\
 & Joint (2-population) SFS & -- \\
\midrule

% Admixture / F-statistics
\multirow{3}{*}{Admixture} & Patterson's $F_2$ & \citep{patterson2012} \\
 & Patterson's $F_3$ & \citep{patterson2012} \\
 & Patterson's $D$ (ABBA-BABA) & \citep{patterson2012} \\
\midrule

% Dimensionality reduction
\multirow{3}{*}{Dim.\ reduction} & PCA & \citep{patterson2006} \\
 & Randomized PCA & \citep{halko2011} \\
 & PCoA (classical MDS) & -- \\
\midrule

% Relatedness
\multirow{2}{*}{Relatedness} & GRM & \citep{yang2011} \\
 & IBS & -- \\
\midrule

% Distance distribution
\multirow{4}{*}{Distance dist.} & Pairwise Hamming distance & -- \\
 & Variance of distances & \citep{schrider2018} \\
 & Skewness of distances & \citep{schrider2018} \\
 & Excess kurtosis of distances & \citep{schrider2018} \\

\end{longtable}

\end{document}
